\section{Trộn lại ở mỗi bước}
\label{sec:ch06-tron-lai}

\index{attention!ba phương trình}%
Cơ chế gọn hơn tôi tưởng trước khi đọc kỹ. Ba phương trình, và phương trình
thứ hai làm gần hết việc.

Encoder đọc \tn{câu nguồn}{source sentence} và để lại một vector cho mỗi vị
trí. Bài gọi chúng là annotation, và tôi giữ nguyên chữ ấy: \(\vect{h}_1\) tới
\(\vect{h}_{T}\), \(\vect{h}_j\) ứng với từ nguồn thứ \(j\). Chương trước cũng
tính đúng chừng ấy vector rồi bỏ đi tất cả trừ cái cuối. Ở đây không cái nào bị
bỏ.

Decoder ở bước \(i\) tự dựng lấy vector ngữ cảnh của riêng nó:

\begin{equation}
  \vect{c}_i = \sum_{j=1}^{T} \alpha_{ij} \vect{h}_j,
  \label{eq:ch06-tong-co-trong-so}
\end{equation}

đúng phương trình (5) của bài. Đây là toàn bộ ý tưởng. Chỉ số \(i\) ở vế trái
là chỗ đáng nhìn lâu: chương trước có một \(\vect{v}\) duy nhất cho cả câu, ở
đây mỗi bước đích có một \(\vect{c}_i\) riêng, dựng từ cùng một bộ annotation
nhưng với bộ trọng số khác.

Các trọng số ấy là một phân phối trên các vị trí nguồn:

\begin{equation}
  \alpha_{ij} = \frac{\exp(e_{ij})}{\sum_{k=1}^{T} \exp(e_{ik})},
  \qquad
  e_{ij} = a(\vect{s}_{i-1}, \vect{h}_j),
  \label{eq:ch06-trong-so}
\end{equation}

phương trình (6). \(e_{ij}\) là \tn{điểm đồng chỉnh}{alignment score} giữa bước
đích \(i\) và vị trí nguồn \(j\); softmax biến các điểm ấy thành
\tn{trọng số đồng chỉnh}{alignment weight} cộng lại bằng 1. Còn \(a\) là một
mạng nhỏ, huấn luyện chung với mọi thứ còn lại.

\subsection*{Điểm chấm bằng trạng thái trước bước, không phải sau}

\index{attention!thứ tự trong một bước decoder}%
Chi tiết dễ đọc lướt nhất trong cả mục này nằm ở chỉ số của \(\vect{s}\) trong
\eqref{eq:ch06-trong-so}. Không phải \(\vect{s}_i\) mà là \(\vect{s}_{i-1}\).
Mục 3.1 của bài nói rõ ra bằng chữ: điểm được chấm dựa trên trạng thái
\(\vect{s}_{i-1}\), tức trạng thái ngay trước lúc sinh ra \(y_i\)
(\enquote{just before emitting \(y_i\)}).

Nên một bước decoder chạy theo thứ tự này, và thứ tự mới là thứ đáng nhớ chứ
không phải công thức:

\begin{enumerate}
  \item lấy trạng thái hiện có \(\vect{s}_{i-1}\), chấm điểm nó với từng
    annotation, softmax ra \(\alpha_{ij}\);
  \item trộn annotation theo trọng số ấy thành \(\vect{c}_i\);
  \item bây giờ mới bước: \(\vect{s}_i = f(\vect{s}_{i-1}, \vect{y}_{i-1},
    \vect{c}_i)\), tức phương trình (4);
  \item đọc phân phối từ tiếp theo ra khỏi \(\vect{s}_i\).
\end{enumerate}

Nghĩa là mô hình quyết định nhìn vào đâu \emph{trước} khi bước, dựa trên chỗ nó
đang đứng. Luong và cộng sự làm ngược lại một năm sau, bước trước rồi mới chấm
điểm bằng trạng thái mới, và họ liệt kê chỗ khác nhau ấy như một trong ba khác
biệt chính. Hộp cầu nối ở mục~\ref{sec:ch06-cau-noi} là chỗ nói chuyện đó.

Trong repo companion, bốn bước trên là bốn dòng:

\begin{minted}{python}
def step(self, word, state, annotations, mask, projected):
    alpha = self.attention.weights(state[1], projected, mask)
    context = self.attention.context(alpha, annotations)
    merged = torch.cat([word, context], dim=-1)
    return self.decoder.step(merged, state), alpha
\end{minted}

\code{state} đi vào \code{weights} trước khi đi vào \code{decoder.step}, và đó
là toàn bộ chỗ khác nhau giữa bản này với bản của Luong.

\subsection*{Mô hình chấm điểm là một mạng một tầng ẩn}

\index{attention!hàm chấm điểm của Bahdanau}%
Phụ lục A.1.2 của bài chọn dạng cụ thể cho \(a\):

\begin{equation}
  a(\vect{s}_{i-1}, \vect{h}_j)
  = \vect{v}_a\transpose \tanh\!\left(
      \mat{W}_a \vect{s}_{i-1} + \mat{U}_a \vect{h}_j
    \right).
  \label{eq:ch06-ham-cham-diem}
\end{equation}

Một tầng ẩn, tanh, rồi chiếu xuống một số thực. Bài giải thích lựa chọn bằng
chi phí: mô hình này phải chạy \(T_x \times T_y\) lần cho một cặp câu, nên nó
phải rẻ. Con số \(T_x \times T_y\) ấy quay lại ở mục~\ref{sec:ch06-con-lai-gi}
và nó không nhỏ như bài nghĩ.

\index{attention!tính trước U\_a h\_j}%
Có đúng một phép tối ưu mà bài viết ra, và nó nằm gọn trong cấu trúc của
\eqref{eq:ch06-ham-cham-diem}. Số hạng \(\mat{U}_a \vect{h}_j\) không chứa
\(i\). Nó là cùng một giá trị ở mọi bước đích, nên tính một lần cho cả câu
thay vì \(T_y\) lần. Nguyên văn: \enquote{Since \(U_a h_j\) does not depend on
\(i\), we can pre-compute it in advance to minimize the computational cost}.

\begin{minted}{python}
def precompute(self, annotations: torch.Tensor) -> torch.Tensor:
    return self.u_a(annotations)
\end{minted}

\begin{minted}{python}
def weights(
    self,
    state: torch.Tensor,
    projected: torch.Tensor,
    mask: torch.Tensor,
) -> torch.Tensor:
    hidden = torch.tanh(self.w_a(state) + projected)
    scores = self.v_a(hidden).squeeze(-1)
    scores = scores.masked_fill(~mask, float("-inf"))
    return torch.softmax(scores, dim=0)
\end{minted}

\code{projected} là kết quả của \code{precompute}, truyền vào từ ngoài. Cấu
trúc cộng trong \eqref{eq:ch06-ham-cham-diem} là thứ cho phép chuyện đó: nếu
hàm chấm điểm nhân \(\vect{s}\) với \(\vect{h}\) thay vì cộng hai phép chiếu
thì không tách ra được.

\index{attention!mask trên vị trí đệm}%
Dòng \code{masked\_fill} là dòng dễ quên nhất trong cả module, và nó hỏng theo
đúng kiểu mà chương~\ref{ch:encoder-decoder} đã gặp một lần: im lặng. Một câu
ngắn nằm chung batch với câu dài thì các cột thừa là token đệm. Không có dòng
ấy, softmax chạy trên cả các cột đệm, nên câu ngắn tiêu một phần ngân sách chú
ý vào chỗ trống. \tn{Hàm mất mát}{loss function} không kêu gì cả, và cái duy
nhất bạn thấy là vài
điểm chính xác biến mất. Repo có test riêng khẳng định trọng số ở vị trí đệm
đúng bằng 0.

\subsection*{Trạng thái khởi tạo lấy từ chiều ngược}

\index{attention!trạng thái khởi tạo của decoder}%
Phụ lục A.2.2 đặt \(\vect{s}_0 = \tanh(\mat{W}_s \overleftarrow{\vect{h}}_1)\).
Không phải trạng thái cuối của chiều xuôi, mà là trạng thái của chiều ngược tại
vị trí nguồn \emph{1}. Nghe ngược đời cho tới khi nhớ rằng chiều ngược đọc từ
phải sang trái: trạng thái của nó ở vị trí 1 chính là trạng thái sau khi nó đã
đọc hết câu. Chương trước dùng trạng thái cuối của chiều xuôi cho đúng việc ấy,
và khác biệt chỉ là đọc từ đầu nào.

\begin{minted}{python}
def encode(self, source: torch.Tensor):
    embedded = self.source_embedding(source)
    annotations, mask = self.encoder(source, embedded)
    projected = self.attention.precompute(annotations)
    h_0 = torch.tanh(self.bridge(self.encoder.summary(annotations)))
    state = (torch.zeros_like(h_0), h_0)
    return annotations, mask, projected, state
\end{minted}

Hình~\ref{fig:ch06-kien-truc} ghép cả mục này lại thành một bức.

\begin{figure}[htbp]
  \centering
  % Kiến trúc của Bahdanau 2015, vẽ theo hình 1 của bài.
% Encoder hai chiều cho mỗi vị trí nguồn một annotation; decoder ở bước i chấm
% điểm mọi annotation bằng s_{i-1}, rồi lấy tổng có trọng số làm c_i.
% Mono-safe: chỉ nét liền, nét đứt, độ dày nét và mức xám.
\begin{tikzpicture}[
  >= Stealth,
  buoc/.style   = {draw, rectangle, minimum width = 7mm, minimum height = 5mm,
                   inner sep = 0pt},
  gop/.style    = {draw, rectangle, minimum width = 9mm, minimum height = 15mm,
                   inner sep = 0pt, black!55, densely dashed},
  mui/.style    = {->, line width = 0.6pt, black!75},
  trongso/.style = {->, line width = 0.5pt, black!55},
  ngucanh/.style = {->, line width = 1.6pt, black},
  nho/.style    = {font = \scriptsize},
  ghichu/.style = {font = \footnotesize, black!60},
]

  % --- bốn vị trí nguồn, mỗi vị trí một cột
  \foreach \j/\lab in {0/1, 1/2, 2/3, 3/{T_x}} {
    \pgfmathsetmacro\x{\j*1.6}
    \node[buoc] (f\j) at (\x, 0.45) {};
    \node[buoc] (b\j) at (\x, -0.45) {};
    \node[gop]  (h\j) at (\x, 0) {};
    \draw[mui] (\x, -1.5) -- (\x, -0.8);
    \node[nho, anchor = north] at (\x, -1.55) {\(\vect{x}_{\lab}\)};
  }

  % --- chiều xuôi đọc trái sang phải, chiều ngược đọc phải sang trái
  \foreach \j [evaluate = \j as \k using int(\j+1)] in {0, 1, 2} {
    \draw[mui] (f\j) -- (f\k);
    \draw[mui] (b\k) -- (b\j);
  }

  \node[ghichu, anchor = west] at (5.45, 0.45) {chiều xuôi};
  \node[ghichu, anchor = west] at (5.45, -0.45) {chiều ngược};
  \node[nho, anchor = north] at (2.4, -2.15)
    {\(\vect{h}_j = [\overrightarrow{\vect{h}}_j ; \overleftarrow{\vect{h}}_j]\)};

  % --- trọng số đồng chỉnh, hội tụ về một nút tổng
  \node[draw, circle, inner sep = 1pt] (tong) at (2.4, 1.9) {\(\Sigma\)};
  \foreach \j in {0, 1, 2, 3} {
    \draw[trongso] (h\j) -- (tong);
  }
  \node[nho, anchor = east] at (0.55, 1.35) {\(\alpha_{i1}\)};
  \node[nho, anchor = west] at (4.3, 1.35) {\(\alpha_{iT_x}\)};

  % --- vector ngữ cảnh của riêng bước i
  \draw[ngucanh] (tong) -- (5.4, 1.9) node[midway, above, nho] {\(\vect{c}_i\)};

  % --- hai bước decoder
  \node[buoc, minimum width = 9mm] (truoc) at (6.05, 2.75) {\(\vect{s}_{i-1}\)};
  \node[buoc, minimum width = 9mm] (sau)   at (7.85, 2.75) {\(\vect{s}_i\)};
  \draw[mui] (truoc) -- (sau);
  \draw[mui] (5.4, 1.9) -- (7.85, 1.9) -- (sau);
  \draw[mui] (sau) -- (7.85, 3.6);
  \node[nho, anchor = south] at (7.85, 3.65) {\(\vect{y}_i\)};

  % --- điểm chấm bằng trạng thái TRƯỚC bước, không phải sau.
  % Không gắn nhãn chữ cho mũi tên này: khoảng giữa nút tổng và hộp
  % \(\vect{s}_{i-1}\) đã có mũi tên \(\vect{c}_i\) chạy qua, và mọi chỗ đặt
  % chữ thử qua đều đè lên một trong hai. Chú thích hình nói thay.
  \draw[->, line width = 0.5pt, black!75, densely dashed]
    (truoc.north west) to[out = 150, in = 55] (tong);

\end{tikzpicture}

  \caption{Kiến trúc của bài, vẽ theo hình 1 của nó. So với
    hình~\ref{fig:ch05-encoder-decoder} thì chỗ nối giữa hai nửa đã đổi hẳn
    hình dạng: mũi tên đậm không còn xuất phát từ một ô duy nhất ở cuối
    encoder, nó xuất phát từ một tổng có trọng số lấy trên \emph{mọi} vị trí
    nguồn, và tổng ấy được tính lại ở từng bước đích. Mũi tên đứt là chi tiết
    dễ mất nhất: điểm đồng chỉnh chấm bằng \(\vect{s}_{i-1}\), tức trạng thái
    trước bước, nên vector ngữ cảnh của bước \(i\) đã được chọn xong trước khi
    bước \(i\) chạy.}
  \label{fig:ch06-kien-truc}
\end{figure}

\subsection*{Hai chỗ repo lệch khỏi bài, và một chỗ bài cho phép lệch}

\index{attention!chỗ repo lệch khỏi bài}%
Bài dùng đơn vị có cổng của Cho, tức gated recurrent unit, viết tắt là GRU, ở
cả encoder lẫn decoder. Repo dùng ô nhớ LSTM có \tn{cổng quên}{forget gate} của
chương~\ref{ch:encoder-decoder}. Đây không phải chỗ tôi tự tiện: phụ lục A.1.1
của chính bài viết rằng \enquote{it is therefore possible to use LSTM units
instead of the gated hidden unit described here, as was done in a similar
context by Sutskever et al. (2014)}. Và lý do tôi chọn thế mạnh hơn là chuyện
được phép: chương này sẽ đặt bảng của nó cạnh bảng chương trước, nên hai bên
phải khác nhau đúng một thứ. Đổi cả ô nhớ thì không còn quy được chênh lệch cho
attention nữa.

Chỗ lệch thứ hai: bài đưa trạng thái qua một tầng maxout rồi mới tới softmax,
repo đọc thẳng softmax từ \(\vect{s}_i\). Cùng lý do, cộng thêm chuyện
chương~\ref{ch:encoder-decoder} cũng đọc thẳng.

Còn một chỗ lệch mà tôi \emph{không} cho phép mình: kích thước. Bài chạy
\(n = 1000\) đơn vị ẩn, embedding 620 chiều, tầng maxout 500, mô hình chấm điểm
1000 đơn vị, huấn luyện năm ngày. Repo chạy vài chục đơn vị trong vài chục
giây. Mọi con số ở các mục sau phải đọc trong khung ấy, và
mục~\ref{sec:ch05-corpus} đã liệt kê corpus này không mô hình hóa những gì.

\subsection*{Một tính chất của lúc khởi tạo, để dành cho mục sau}

\index{attention!v\_a khởi tạo bằng 0}%
Phụ lục B.1 khởi tạo \(\mat{W}_a\) và \(\mat{U}_a\) từ phân phối chuẩn với độ
lệch chuẩn 0.001, và khởi tạo \(\vect{v}_a\) \emph{đúng bằng 0}.

Đặt vào \eqref{eq:ch06-ham-cham-diem} thì mọi \(e_{ij}\) bằng 0 ở batch đầu
tiên, nên softmax trả về phân phối đều: \(\alpha_{ij} = 1/T\) với mọi \(j\).
Mô hình khởi động bằng cách lấy trung bình cộng tất cả annotation, và phải học
lấy mới biết ưu tiên cái nào.

Tôi ghi chi tiết này ra thay vì dọn nó đi vì mục sau đo gradient tại đúng thời
điểm ấy, và một phân phối đều là điểm xuất phát sạch nhất có thể có cho phép đo
đó.
